\section{Maximum rank: one transposition for each tree edge}
\begin{theorem}\label{thm:maxrank}
Let $n\geq7$ and let $G\leq S_n$ have a C-group marking of rank
$n-1$. Then $G=S_n$, and its distinguished generators are exactly
the transpositions on the edges of a tree on the $n$ letters.
Conversely every such tree gives a C-group marking of rank $n-1$.
Its Coxeter diagram is the line graph of the tree, and its coset system
is a regular hypertope with full type-preserving group $S_n$.
\end{theorem}
\begin{proof}
The marking is an independent family of length $n-1$, so the equality
clause of Whiston's theorem gives $G=S_n$. Retain the omission
alternatives recorded in [CC, proof of Theorem 2.1]: each maximal
standard parabolic is intransitive, transitive imprimitive with blocks
of size two, or $A_n$. This is an explicit outside input; we do not
infer it from the rank-$n-2$ source result being recovered.

In the imprimitive case write $n=2m$, so $m\geq4$. The binary-module
proof in [U] gives an irredundant-length bound $2m-4$ if the block
image is a proper subgroup of $S_m$, and $m+1$ if it is $S_m$.
Both are strictly smaller than $2m-2$, even for $m=4$. But our
parabolic has $n-2=2m-2$ independent generators. The $A_n$ case is
excluded by the elementary intersection argument in [U]: the omitted
involution must centralize every other generator, whereas the
centralizer of $A_n$ in $S_n$ is trivial. Hence every maximal
parabolic $G_i$ is intransitive.

For every $i$, choose an $i$-edge joining two distinct $G_i$-orbits
in the permutation graph. These $n-1$ edges form a forest: in a
cycle the selected $i$-edge could be bypassed using other labels,
contradicting its orbit choice. Thus they form a spanning tree $T$.
An extra $i$-edge would, with its path in $T$, form a cycle. Every
selected edge on that path must have label $i$, by the same bypass
argument. Since only one selected $i$-edge exists, that path is the
selected edge itself, not a new edge. There are no extra edges,
and all distinguished permutations are the tree transpositions.

Conversely the transpositions on a connected tree generate $S_n$.
For any edge subset $F$, they generate the product of symmetric groups
on the connected components of $(V(T),F)$. Two vertices are in a common
component of both $F$ and $E$ exactly when their unique tree path lies
in $F\cap E$. Intersecting the two Young subgroups therefore gives
precisely the subgroup for $F\cap E$, proving the full intersection
property. Two different edge transpositions have product order three
when their edges meet and order two otherwise, giving the asserted
line graph.

For the geometric claim root $T$ at any vertex. An object of the type
corresponding to an edge $e$ is a subset of size equal to the component
of $T-e$ away from the root, because its maximal parabolic is the
stabilizer of that subset. Two such objects are incident precisely when
they have the prescribed disjointness or containment. For a flag of
these subsets, subtract the maximal selected children from each selected
set and take the complement of the maximal selected sets for the root
region. The resulting disjoint regions have exactly the cardinalities
of the corresponding standard regions of $T$. Independent bijections
on the regions combine to a permutation realizing the whole flag.
Thus the system is flag-transitive. The elementary coset argument
proved in Section~\ref{sec:regularity} gives thinness, residual
connectedness and regularity, including the full automorphism group.
\end{proof}
The uniform tree-presentation theorem of [U] now recovers
[FL, Proposition 3.1]: use the Coxeter relations of the line graph
and the triangle relations $(s_es_fs_es_g)^2=1$.
That proof is valid for every finite tree, even when $n<7$; the
classification theorem above retains its stated range $n\geq7$.
